
\begin{abstract}
  华北山区乡村的耕地类型多、单块面积差异大，种植安排还同时受到季节适宜性、重茬和豆类轮作等约束。本文以54个地块、41种作物及2023年种植统计为基础，研究2024至2030年的种植面积配置，并把经济收益、下行风险和田间管理纳入同一套分析框架。

  \textbf{针对问题一，}首先整理地块、作物和产销参数，统一相邻种植季与三年豆类覆盖的判定标准。其次，以逐年逐季逐地块种植面积为决策变量，建立\textbf{确定性多期混合整数线性规划模型}，分别处理超额产量滞销和半价销售两种模式。然后，统一每地块季次最多3种作物的管理约束，半价模式另施加10\%最小面积占比，并以规则轮作方案作为基线。最后，使用HiGHS求解并独立核验收益和硬约束。正式滞销方案的七年累计净收益为\textbf{4041.72万元}，MIP Gap为\textbf{2.126\%}；半价销售方案为\textbf{6351.02万元}，MIP Gap为\textbf{0.739\%}，两份方案的硬约束违反数均为0。

  \textbf{针对问题二，}首先依据题目给出的变化区间构造逐年复利的中心趋势。其次，采用\textbf{期望趋势MILP与样本外蒙特卡洛评估模型}，先求唯一面积方案，再固定该方案进入200个随机情景，避免复制大规模情景决策变量。然后，与2023基期静态模型在共同情景中配对比较。最后，使用五个随机种子、三角分布和端点情景检查稳定性。半价模式下主方案平均七年净收益为\textbf{6897.19万元}，比基线高\textbf{5.68万元}，最差5\%情景平均收益为\textbf{6317.80万元}。正式测试及稳健性检查中均未出现模拟亏损，但下尾指标并非始终优于基线。

  \textbf{针对问题三，}首先按作物类别与用途构造替代和互补响应，并明确其属于模拟假设。其次，在问题二框架上建立\textbf{关系调整趋势MILP与相关蒙特卡洛模型}，设置弱、中、强三档关系强度。然后，直接在最终销量、亩产量和价格情景上检验经验相关矩阵。最后，从收益和41维作物累计面积两个层面判断结构变化。中档方案半价模式平均收益为\textbf{6887.03万元}，比关闭关系结构的基线高\textbf{12.14万元}；三档方案两两收益差为\textbf{0.061\%至0.186\%}，面积结构相似度为\textbf{94.73\%至96.63\%}，说明宏观种植结构稳定，差异主要来自地块间的微观调度。

  本文的主要工作有三个方面。(1) 将产销分段收益、轮作和管理便利统一写入可核验的多期MILP，(2) 将方案求解与随机风险评估解耦，在保留完整土地约束的同时控制计算规模，(3) 对问题三采用宏观结构与微观调度双层归因，避免用低地块重合度误判市场结构跳变。所得结论只适用于题目给定区间和本文明确列出的模拟假设。

  \keywords{混合整数线性规划；种植策略；蒙特卡洛模拟；下尾风险；敏感性分析}
\end{abstract}
